\section{Síntesis y ampliación}

Esta sección condensa EP112 en cinco conclusiones. Primera: las empresas de modelos empiezan a diferenciarse, y lo general, Coding+Agentic, la interacción multimodal y el posicionamiento en el ecosistema se vuelven rutas distintas. Segunda: el punto de acceso para el consumidor converge, y la suite horizontal de ChatGPT y la integración vertical de Gemini representan dos formas de absorción. Tercera: tanto la inteligencia como el producto importan, y las barreras de producto empiezan a valorarse de nuevo. Cuarta: los productos de IA se parecen a la minería, y la ventana del Magic Moment es cada vez más corta. Quinta: la experiencia L4 ya apareció en Deep Research y Claude Code, y el mayor beneficio actual sigue estando en language/code.

\begin{importantbox}{EP112 dentro de la serie de IA de Zhang Xiaojun (张小珺)}
EP127 y EP136 son informes trimestrales más amplios sobre los modelos grandes; EP115/EP113 tratan sobre Agent y la ruta técnica de K2. EP112, en cambio, regresa la mirada al producto y al mercado: cómo se diferencian las empresas de modelos, cómo los productos de IA compiten por su ventana y cómo convergen las rutas de plataforma de Google/OpenAI/Anthropic.
\end{importantbox}

\subsection{Takeaways clave}

Las secciones anteriores abordaron por separado la diferenciación de las empresas, la convergencia de las plataformas, la ventana de producto y el relato de la burbuja; esta sección los condensa en juicios aplicables. Para el inversionista, lo importante no es seguir cada lanzamiento de un modelo nuevo, sino determinar si la capacidad se convierte en punto de acceso e ingresos; para el emprendedor, lo importante no es demostrar que el modelo puede hacerlo, sino demostrar que puede retener a los usuarios una vez cerrada la ventana.

\begin{knowledgebox}{Lista de juicios de este episodio}
\begin{tabularx}{\textwidth}{p{0.24\textwidth}p{0.34\textwidth}X}
\toprule
Juicio & Significado & Indicadores a observar \\
\midrule
Diferenciación de las empresas de modelos & Cada empresa elige una pila de capacidades distinta & Rankings de coding, agentic, multimodalidad y uso general. \\
El producto vuelve a importar & La inteligencia debe convertirse en punto de acceso y en posicionamiento en la mente del usuario & DAU, retención, pago, elección predeterminada de marca. \\
La ventana se acorta & El Magic Moment se copia más rápido & Velocidad de respuesta de la competencia, absorción en funciones nativas del modelo. \\
Google se revalúa & Activos antiguos pueden recombinarse en la IA & Integración de Search, Ads, Gemini y Workspace. \\
La burbuja depende de que se cumpla lo prometido & El ánimo del mercado no es el único criterio & Ingresos, CapEx, avance de los modelos y productos L4. \\
\bottomrule
\end{tabularx}
\end{knowledgebox}

\begin{knowledgebox}{Marco de decisión para invertir/emprender}
\begin{tabularx}{\textwidth}{p{0.22\textwidth}p{0.34\textwidth}X}
\toprule
Pregunta & Qué observa el inversionista & Qué debe responder el emprendedor \\
\midrule
¿La capacidad es escasa? & ¿Cuánto tardarán las empresas de modelos en alcanzarla? & ¿Cuánto dura mi ventana de experiencia? \\
¿El punto de acceso es claro? & ¿Los usuarios lo abren con frecuencia? & ¿Puedo convertirme en el flujo de trabajo predeterminado? \\
¿Los datos regresan al producto? & ¿Mejora mientras más se usa? & ¿Puedo acumular datos o preferencias propios? \\
¿El modelo de negocio se sostiene? & ¿Los ingresos cubren el costo de inferencia? & ¿Cómo convierto el Magic Moment en una razón para pagar? \\
¿Lo harán los gigantes? & ¿La suite completa ya lo incluye de forma natural? & ¿Puedo evitar la plataforma o apoyarme en ella? \\
\bottomrule
\end{tabularx}
\end{knowledgebox}

\begin{enumerate}
\item La diferenciación de las empresas de modelos no es un debilitamiento, sino una elección de ruta: lo general, Coding+Agentic y la interacción multimodal compiten por posiciones dominantes distintas.
\item La suite horizontal y la integración vertical seguirán reduciendo el espacio de supervivencia de los productos de IA independientes.
\item Las barreras de producto incluyen el punto de acceso predeterminado, el posicionamiento de marca, el retorno de datos y la monetización, no solo la capacidad del modelo.
\item La Magic Window de los productos de IA se está acortando; el emprendedor debe convertir rápidamente el asombro en una barrera.
\item La experiencia L4 más clara sigue estando en language/code, no en la multimodalidad, los modelos del mundo ni la robótica.
\end{enumerate}

\subsection{Preguntas abiertas}

\begin{enumerate}
\item ¿La próxima experiencia L4 aparecerá en derecho, finanzas, ventas, análisis de datos o en el workflow empresarial?
\item ¿Puede la integración vertical de Google contrarrestar la ventaja de punto de acceso horizontal de ChatGPT?
\item ¿Cómo puede una empresa de producto construir una barrera duradera dentro de una Magic Window de 3 a 6 meses?
\item Si se forma una burbuja de la AGI, ¿el punto de ruptura más probable será la desaceleración del avance de los modelos, ingresos por debajo de lo esperado o un rendimiento insuficiente del CapEx?
\item Una vez que las empresas de modelos conviertan sus capacidades en productos, ¿en qué capas conviene más que permanezcan las empresas de aplicaciones independientes: el punto de acceso, la cadena de herramientas, el flujo de trabajo vertical o la red de datos?
\item Cuando se fusionen la búsqueda, la publicidad, la ofimática y Agent, ¿la competencia entre Google y OpenAI se parecerá más a una guerra de buscadores o a una guerra de sistemas operativos?
\end{enumerate}

\begin{warningbox}{Vigencia de los juicios trimestrales}
Los juicios trimestrales caducan por naturaleza. Lo que de verdad vale la pena conservar no es el liderazgo momentáneo de una empresa, sino el marco de análisis: diferenciación, convergencia, ventana de producto, experiencia L4, integración de plataformas e indicadores de burbuja. Este marco puede seguir usándose para observar el próximo trimestre.
\end{warningbox}

\subsection{Lecturas adicionales}

\begin{itemize}
\item Si le interesa la ruta técnica de Agent, puede contrastarse con la entrevista a Yao Shunyu (姚顺雨) en EP115 y la entrevista a Yang Zhilin (杨植麟) sobre K2 en EP113.
\item Si le interesa la adopción en empresas, puede contrastarse con la entrevista a Wu Minghui (吴明辉) en EP116, para entender qué forma toma el Agentic Model con datos privados en el ámbito ToB.
\item Si le interesan los informes trimestrales más amplios sobre los modelos grandes, puede continuar con las notas de los informes trimestrales de EP127 y EP136.
\end{itemize}

\end{document}
